\documentclass[12pt,a4paper]{article}
\usepackage[T2A]{fontenc}
\usepackage[utf8]{inputenc}
\usepackage[russian]{babel}
\usepackage{amsmath,amssymb}
\usepackage[a4paper,margin=22mm]{geometry}
\usepackage[hidelinks]{hyperref}

\title{Сетка по потенциалу точечного источника и стока\\
с изобарическими границами скважин}
\author{Математическая постановка и численная реализация}
\date{}

\begin{document}
\maketitle

\section{Область и смысл заданного радиуса}

Рассматривается прямоугольная ячейка пятиточечной системы
$0\le x\le L_x$, $0\le y\le L_y$. Точечный источник расположен в
$z_I=0$, точечный сток --- в $z_P=L_x+\mathrm iL_y$. Физические круговые
границы в аналитическую задачу не вводятся. Заданные числа $r_I$ и $r_P$
служат для выбора двух изобар, которые заменяют границы скважин.

Изобара выбирается со стороны пласта: она касается окружности заданного
радиуса хотя бы в одной точке и нигде не проходит ближе к точечной
особенности. Поэтому расчётная область не заходит внутрь номинальной
окружности.

\section{Периодический комплексный потенциал}

Пятиточечная система зеркально продолжается на плоскость. Источники находятся
в узлах $(2mL_x,2nL_y)$, а стоки --- в узлах
$((2m+1)L_x,(2n+1)L_y)$. Такое продолжение обеспечивает нулевую нормальную
скорость на четырёх прямых сторонах основной ячейки.

Для устойчивого суммирования положим
\[
 a=\min(L_x,L_y),\qquad b=\max(L_x,L_y),\qquad
 q=\exp(-\pi b/a).
\]
Если $L_x>L_y$, координаты перед вычислением меняются местами. В преобразованной
координате $z=x'+\mathrm iy'$ сток расположен в $a+\mathrm ib$. Введём
\[
 \Theta(\zeta;q)=\sum_{n=0}^{\infty}(-1)^n
 q^{(n+1/2)^2}\sin((2n+1)\zeta).
\]
Комплексный потенциал равен
\begin{equation}
 F(z)=log\Theta\!\left(\frac{\pi z}{2a};q\right)
 -\log\Theta\!\left(\frac{\pi(z-a-\mathrm ib)}{2a};q\right)
 +\frac{\mathrm i\pi z}{2a}.
 \label{eq:F}
\end{equation}
Здесь
\[
 \varphi=\operatorname{Re}F,\qquad
 \psi=\operatorname{Im}F.
\]
Изобары задаются уравнением $\varphi=\mathrm{const}$, линии тока ---
$\psi=\mathrm{const}$. Производная
\begin{equation}
 F'(z)=\frac{\pi}{2a}\left[
 \frac{\Theta'(\pi z/2a)}{\Theta(\pi z/2a)}-
 \frac{\Theta'(\pi(z-a-\mathrm ib)/2a)}
      {\Theta(\pi(z-a-\mathrm ib)/2a)}+\mathrm i
 \right]
 \label{eq:velocity}
\end{equation}
даёт градиент потенциала:
$\nabla\varphi=(\operatorname{Re}F',-\operatorname{Im}F')$.

\section{Выбор изобарических границ}

На номинальной четверти окружности нагнетателя вычисляется
$\varphi_I(\theta)=\varphi(r_I\cos\theta,r_I\sin\theta)$,
$0\le\theta\le\pi/2$. Уровень эффективной границы выбирается как
\begin{equation}
 \varphi_I^*=\max_{0\le\theta\le\pi/2}\varphi_I(\theta).
 \label{eq:injector-level}
\end{equation}
Вдоль каждого луча решается скалярное уравнение
\[
 \varphi(\rho_I(\theta)\cos\theta,
         \rho_I(\theta)\sin\theta)=\varphi_I^*.
\]
Из монотонности потенциала около источника следует
$\rho_I(\theta)\ge r_I$, а хотя бы для одного $\theta$ выполняется равенство.

Для добывающей скважины точка на луче записывается как
\[
 (x,y)=(L_x-\rho_P\sin\theta,L_y-\rho_P\cos\theta).
\]
Поскольку потенциал растёт при приближении к стоку, выбирается
\begin{equation}
 \varphi_P^*=\min_{0\le\theta\le\pi/2}\varphi_P(r_P,\theta),
 \label{eq:producer-level}
\end{equation}
и затем решается уравнение $\varphi=\varphi_P^*$. В результате также
$\rho_P(\theta)\ge r_P$ и $\min\rho_P=r_P$.

\section{Линии тока без дополнительного сгущения у границ}

На найденной изобаре нагнетателя вычисляется непрерывная ветвь $\psi$.
Диапазон между двумя непроницаемыми внешними путями делится равномерно:
\begin{equation}
 \psi_j=\psi_{min}+\frac{j}{N_\psi}
 (\psi_{max}-\psi_{min}),\qquad j=0,\ldots,N_\psi.
 \label{eq:psi-levels}
\end{equation}
Никакого дополнительного смещения уровней к внешним границам не выполняется.
Число линий около каждого участка определяется самим полем и общим числом
$N_\psi$.

Две крайние линии тока совпадают с ломаными непроницаемыми границами. Их узлы
находятся одномерным обращением $\varphi$ вдоль этих ломаных. Потенциалы двух
внешних углов обязательно включаются в набор уровней, чтобы сетка не срезала
углы прямоугольника.

\section{Естественное логарифмическое сгущение}

Уровни потенциала задаются равномерно:
\begin{equation}
 \varphi_i=\varphi_I^*+\frac{i}{N_\varphi}
 (\varphi_P^*-\varphi_I^*),\qquad i=0,\ldots,N_\varphi.
 \label{eq:phi-levels}
\end{equation}
Около точечного источника
\[
 \varphi=C_I+\log r+o(1).
\]
Следовательно,
\[
 r_i\simeq\exp(\varphi_i-C_I),\qquad
 \frac{r_{i+1}}{r_i}\simeq
 \exp(\varphi_{i+1}-\varphi_i)=\mathrm{const}.
\]
Таким образом, равномерный шаг по $\varphi$ автоматически создаёт
геометрическую прогрессию физических расстояний. Это и есть естественное
логарифмическое сгущение; отдельная мониторная функция около скважин не нужна.
Та же конструкция действует около точечного стока.

\section{Трассировка и сборка сетки}

Каждая внутренняя линия $\psi=\psi_j$ интегрируется с использованием
$\varphi$ как независимой переменной:
\begin{equation}
 \frac{d\mathbf x}{d\varphi}
 =\frac{\nabla\varphi}{|\nabla\varphi|^2}.
 \label{eq:trace}
\end{equation}
Начальная точка находится на изобаре $\varphi_I^*$, а решение выдаётся на
общих уровнях~\eqref{eq:phi-levels} вплоть до $\varphi_P^*$. В программе
используется адаптивная схема Рунге--Кутты. Конечная точка дополнительно
корректируется однимерным шагом Ньютона вдоль $\nabla\varphi$, но не
проецируется на окружность.

Итоговые узлы имеют вид
\[
 \mathbf X_{ij}=\mathbf x(\varphi_i,\psi_j).
\]
Поскольку $F$ аналитична, линии $\varphi=\mathrm{const}$ и
$\psi=\mathrm{const}$ ортогональны вне точечных особенностей и точек
стагнации.

\section{Проверки}

Перед передачей сетки конечно-объёмному решателю проверяются:
\begin{enumerate}
 \item положительность площадей всех ячеек;
 \item минимальный угол и отношение сторон;
 \item принадлежность первого и последнего рядов выбранным изобарам;
 \item условия $\min\rho_I=r_I$, $\min\rho_P=r_P$ с численным допуском;
 \item отсутствие нормальной скорости на прямых внешних границах;
 \item дискретная ортогональность двух семейств координатных линий;
 \item включение обоих внешних углов в границу сетки.
\end{enumerate}

Алгоритм предназначен для геометрии однородной изотропной ячейки. При
неоднородной или анизотропной проницаемости построенные линии уже не совпадают
с фактическими линиями тока расчётной модели.

\end{document}
